%% ---------------------------------------------------------------------
\section{DIS 的扭度 4：从矩到 \texorpdfstring{$x$}{x} 空间}\label{sec:power-twist4}

单举 DIS 是幂次修正唯一有完整算符结构的强子过程。对偶数 $n$ 的结构函数矩，OPE 给出
\begin{equation}\label{eq:power-moments-tw4}
M_2^{(n)}(Q^2)\equiv\int_0^1\!\dd x\,x^{n-2}F_2(x,Q^2)
=\sum_iC^{(n)}_{2,i}A^{(n)}_i+\frac1{Q^2}\sum_jC^{(n),\tau=4}_{2,j}A^{(n),\tau=4}_j+O\Big(\frac{1}{Q^4}\Big),
\end{equation}
其中第一项是扭度 2 的贡献（\cref{thm:dis-moments}），第二项来自扭度 $\tau=4$ 的规范不变定域算符（夸克—胶子算符与四夸克算符），其约化矩阵元 $A^{(n),\tau=4}_j$ 是非微扰参数，系数 $C^{(n),\tau=4}_{2,j}$ 及其反常维度可在微扰论中计算（\cref{prop:dis-twist4}，\tagR，$R_{\mathrm{OPE}}$）。扭度 2 算符的迹项给出另一类 $M^2/Q^2$ 修正，即靶质量修正（\cref{prop:dis-tmc}）；它是运动学的，完全由扭度 2 矩阵元决定，本章所说的“扭度 4”指的是动力学的高扭度。

图像上（\cref{fig:higher-twist}），强子与硬部分之间每多一条横向极化的部分子，贡献就多一个 $1/Q$ 压低。用\cref{lem:dis-twist-classification}中的共线扭度计数：好分量 $\psi_+$、$\bar\psi_+$ 与场强 $G^{+j}$ 的扭度为 1，横向协变导数 $D^j$ 的扭度为 1，$D^+$ 的扭度为 0。于是一条额外的横向胶子给出扭度 3 的夸克—胶子—夸克关联；在非极化 DIS 中它的前向矩阵元没有可用的横向矢量，贡献为零，第一个非零的修正是扭度 4，即 $1/Q^2$。

\paragraph{为什么矩还不够。}
\cref{eq:power-moments-tw4}对每个固定的 $n$ 成立，但给定自旋 $n$ 的扭度 4 算符不是一个而是一族：$n-2$ 个光锥导数 $D^+$ 可以以不同方式分配到各个场上，独立算符的数目随 $n$ 增长。把所有 $n$ 的定域算符重新求和，就得到场位于同一条光锥直线上、但彼此分开的非定域算符，其矩阵元依赖于不止一个动量分数——这正如扭度 2 的定域算符 $\bar\psi\gamma^+(\ii D^+)^{n-1}\psi$ 重新求和后给出部分子分布（\cref{def:pdf-quark}）。$x$ 空间的公式因此是一个同时关于所有 $n$ 的陈述，它在动量空间中的含义是：$1/Q^2$ 修正本身也具有因子化的形式。

\paragraph{扭度 4 关联函数。}
记 $z=(0^+,z^-,\mathbf0_T)$ 为光锥直线上的点。设 $\Phi_1,\dots,\Phi_k$ 取自好分量 $\psi_+$、$\bar\psi_+$ 与横向协变导数 $D_\perp^j$（等价地，在光锥规范 $A^+=0$ 中取自场强 $gG^{+j}$），各场位于光锥直线的不同点上，相邻的场由沿 $n_B$ 的直 Wilson 线 $W[z_a,z_b]$（\cref{def:prelim-wilson}）连成规范不变的组合，且总扭度为 4。对这样的矩阵元 $\langle P|\Phi_1(z_1)W\cdots W\Phi_k(z_k)|P\rangle$ 关于光锥分离作 Fourier 变换，就得到扭度 4 关联函数 $T^{(4)}_i([x])$，其中 $[x]$ 表示独立的光锥动量分数（三个场时有两个，四个场时有三个）。典型代表是 Ellis--Furmanski--Petronzio “横向基”中的 $\bar\psi\gamma^+D_\perp^jD_{\perp j}\psi$ 型夸克—胶子关联与四夸克关联 $\bar\psi\gamma^+t^a\psi\,\bar\psi\gamma^+t^a\psi$；在这个基中，除横向分量外的协变导数都已用运动方程消去\cite{EllisFurmanskiPetronzio1982,EllisFurmanskiPetronzio1983}。在 $A^+=0$ 规范下，这些函数可以表示为夸克、反夸克与胶子的逐次发射与吸收，每个部分子的动量分数都在 $[0,1]$ 内\cite{Jaffe1983}。

\begin{proposition}[DIS 扭度 4 修正的 $x$ 空间形式 \tagA]\label{prop:power-dis-twist4-x}
设 $0<x<1$ 固定，$W^2=(P+q)^2\gg\Lambda^2$。当 $Q\to\infty$ 时，强子靶的非极化结构函数满足
\begin{equation}\label{eq:power-twist4-x}
F_2(x,Q^2)=F_2^{\mathrm{tw2}}(x,Q^2)+\frac{1}{Q^2}\sum_i\int\!\dd x_1\,\dd x_2\,H_i\Big(x_1,x_2,x;\frac{Q^2}{\mu^2},\as(\mu)\Big)\,T^{(4)}_i(x_1,x_2;\mu)+R_4(x,Q),
\end{equation}
其中 $F_2^{\mathrm{tw2}}$ 是扭度 2 的贡献——\cref{thm:disall-factorization}的领头项，加上完全由扭度 2 矩阵元决定的靶质量修正（其 $x$ 空间形式与处方有关，见\cref{prop:dis-tmc}），$T^{(4)}_i$ 是上述规范不变的扭度 4 光锥关联函数（$i$ 遍历一组完备的夸克—胶子与四夸克关联；对四场关联函数，$\int\dd x_1\dd x_2$ 理解为对全部独立动量分数的积分），$H_i$ 是与靶无关、可在微扰论中计算的硬函数，余项 $R_4=O(\Lambda^4/Q^4)$（至多差对数），其系数依赖于 $x$，在 $x\to1$ 与 $x\to0$ 时不均匀。此外：
\begin{enumerate}[label=(\alph*),leftmargin=2em,itemsep=0pt,topsep=2pt]
\item \eqref{eq:power-twist4-x}中 $1/Q^2$ 项的矩 $\int_0^1\dd x\,x^{n-2}(\cdots)$ 等于\eqref{eq:power-moments-tw4}的扭度 4 项，即\cref{prop:dis-twist4}的算符结构；
\item 对核靶，$T^{(4)}_i$ 含有双重散射的部分，按 $A^{1/3}$ 增强（\cref{prop:disall-nuclear}）。
\end{enumerate}
非微扰地位：\cref{asm:intro-bridge}（MA1、MA2、MA4）。
% errtype: E-S
\uses{prop:dis-twist4, prop:dis-tmc, thm:disall-factorization, asm:intro-bridge}
\end{proposition}

\begin{proofsketch}
\begin{enumerate}[label=\arabic*.,leftmargin=1.6em,itemsep=1pt,topsep=2pt]
\item \emph{矩空间。}由\cref{prop:dis-twist4}（以 OPE 为前提而严格），每个偶数矩的 $1/Q^2$ 项是扭度 4 定域算符矩阵元与系数函数之积的和，即\eqref{eq:power-moments-tw4}。
\item \emph{光锥重新求和。}自旋 $n$ 的扭度 4 定域算符是上述非定域光锥算符在各分离点处 Taylor 展开的系数：把 $n-2$ 个 $D^+$ 分配到各个场上，对应于关于各光锥分离的各阶导数。因此 $A^{(n),\tau=4}_j$ 是 $T^{(4)}_i$ 的多重矩 $\int[\dd x]\,x_1^{a_1}x_2^{a_2}\cdots T^{(4)}_i([x])$ 的线性组合，$a_1+a_2+\dots=n-2$。若\eqref{eq:power-twist4-x}成立，则对 $x$ 取矩并交换积分次序即得 (a)，$H_i$ 的矩给出 $C^{(n),\tau=4}_{2,j}$。反之，固定 $n$ 的矩关系只确定 $H_i$ 的有限多个矩，不能推出 $x$ 空间公式：后者需要下面的图形论证。
\item \emph{共线展开。}在割图表示中，把硬部分在部分子横动量 $\mathbf k_T$ 处展开到二阶（\cref{fig:higher-twist}(a)），并加上强子另外经一条或两条共线胶子与硬部分相连的图（(b)、(c)）。用 Ward 恒等式把纵向极化的胶子吸收进 Wilson 线，横动量的一、二阶项与横向胶子场组合成规范不变的横向协变导数，所有贡献都化为“硬函数 $\otimes$ 扭度 4 关联函数”的形式\cite{EllisFurmanskiPetronzio1982,EllisFurmanskiPetronzio1983}。树图层次的硬函数在规范算符基\cite{Jaffe:1982pm}与横向基\cite{EllisFurmanskiPetronzio1983}中都已算出；在光锥规范中所有 $O(1/Q^2)$ 修正都可以用只依赖于纵向动量分数的多部分子分布表示\cite{Jaffe1983}；以光子—部分子前向振幅的完整虚部为出发点、使规范不变性显然的另一种推导见\cite{Qiu:1988dn}。
\item \emph{只被论证的一步。}要使\eqref{eq:power-twist4-x}在动量空间中成立，还须知道：在 $1/Q^2$ 这一阶，主导区域恰好是领头区域展开到二阶、以及额外一两条共线部分子接入硬部分的区域，软区与 Glauber 区不产生新的 $1/Q^2$ 贡献（DIS 只有一个强子，末态被单举地求和，末态相互作用由幺正性相消，参见\cref{thm:disall-factorization}）。这是次领头幂次的区域分析（此步只给出 E-S 型陈述），在文献中是以幂次计数的方式论证的，没有积分范数意义下的估计。
\item \emph{核靶。}核靶中的双重散射使 $T^{(4)}_i$ 获得与核尺度成正比的部分，见\cref{prop:disall-nuclear}。
\end{enumerate}
完整论证见\cite{EllisFurmanskiPetronzio1983}（横向基中的推导及其与 OPE 的对应）与\cite{Qiu:1988dn}。
\end{proofsketch}

\paragraph{地位与缺口}
矩空间的算符结构是 $R_{\mathrm{OPE}}$ 严格的（\cref{prop:dis-twist4}）；$x$ 空间形式\eqref{eq:power-twist4-x}的等级是 \tagA，最弱的一步是上面第 4 步的次领头幂次区域分析（E-S 型），它相当于把\cref{asm:intro-bridge}(Br1) 的“区域完备性”推广到次领头幂次，而本书的编号假设只覆盖领头幂次。领头扭度部分 $F_2^{\mathrm{tw2}}$ 经由\cref{thm:disall-factorization}（\tagA）继承该定理的全部瓶颈——\cref{asm:appN-collins}、\cref{asm:appN-envelope}、\cref{asm:appN-separation}、\cref{asm:appN-cutsum}、\cref{asm:appN-ward-product}、\cref{asm:appN-uv}、\cref{asm:appN-multren}——以及该论证中被明确假定的一步（\cref{asm:appN-brst-nonlocal}，[C]；文献中的位置是\cite[第~11.3.2~节]{Collins2011}）。此外有两点使“扭度 4 修正”本身的含义受到限制。其一，扭度 4 贡献与领头扭度通过重整化子的歧义相混合：扭度 2 系数函数在 $u=1$ 处的红外重整化子歧义与扭度 4 矩阵元的紫外歧义同量级，一般的相消是\cref{conj:power-renormalon-ht}（本章稍后，\tagC）。其二，扭度 2 算符基的选取只在相对 $O(\Lambda^2/Q^2)$ 的精度内有意义，除非用非微扰的方式固定（\cref{prop:power-rearrangement}，本章稍后）。$x$ 空间形式在 $x\to1$ 时的不均匀性见下例；核尺度依赖的控制是\cref{prop:disall-nuclear}的缺口。升级途径：次领头幂次区域分析的积分范数估计（App~N 型的框架推广到 $1/Q^2$），以及\cref{thm:disall-factorization}本身的升级。

\begin{example}[扭度展开在 $x\to1$ 处的增强]\label{ex:power-twist4-largex}
在矩空间中，扭度 4 修正对每个固定的 $n$ 都是 $\Lambda^2/Q^2$ 量级；但它在 $x$ 空间中的大小取决于它随 $n$ 增长的方式。设扭度 2 分布在 $x\to1$ 处的行为是 $q(x)=(1-x)^b$（$b>0$），其 Mellin 矩为
\begin{equation*}
M_n[q]\equiv\int_0^1\!\dd x\,x^{n-1}(1-x)^b=B(n,b+1)=\frac{\Gamma(n)\Gamma(b+1)}{\Gamma(n+b+1)}\;\xrightarrow{\;n\to\infty\;}\;\Gamma(b+1)\,n^{-b-1}.
\end{equation*}
考虑 $x$ 空间中的修正 $\delta q(x)=c\,b\,(1-x)^{b-1}$。它的矩是
\begin{equation*}
M_n[\delta q]=c\,b\,B(n,b)=c\,b\,\frac{\Gamma(n)\Gamma(b)}{\Gamma(n+b)}=c\,(n+b)\,M_n[q],
\end{equation*}
最后一步用了 $\Gamma(b+1)=b\Gamma(b)$ 与 $\Gamma(n+b+1)=(n+b)\Gamma(n+b)$。于是“矩的修正随 $n$ 线性增长”（$\delta M_n\simeq c\,n\,M_n$）与“$x$ 空间的相对修正为 $\delta q/q=c\,b/(1-x)$”是同一件事。取 $c\sim\Lambda^2/Q^2$，扭度展开在大 $x$ 处的实际展开参数是 $\Lambda^2/[Q^2(1-x)]$，当 $1-x\lesssim\Lambda^2/Q^2$ 时展开失效。物理上，这对应于末态不变质量 $W^2\simeq Q^2(1-x)/x$ 变小：大 $x$ 处真正的硬标度是 $Q\sqrt{1-x}$。

这种线性增长确实出现在扭度 4 的“重整化子模型”中。在大 $\beta_0$ 近似下，扭度 2 系数函数在 $u=1$ 处的红外重整化子留数给出一个随 $x$ 变化的形状函数，其非极化 $F_2$ 部分含有 $-4/[1-x]_+$ 与 $-\delta'(1-x)$ 两项\cite[式~(5.33)]{Beneke1999}。前者的 Mellin 矩是 $4S_1(n-1)\sim4\ln n$，后者是 $-(n-1)$（因为 $\int_0^1\dd x\,x^{n-1}\delta'(1-x)=n-1$）；与 $q$ 卷积后，$\delta'(1-\xi)$ 一项给出 $-xq'(x)-q(x)$，在 $x\to1$ 时相对于 $q$ 正是 $b/(1-x)$ 的增强。一般地，与单胶子线相联系的重整化子幂次修正按 $[\Lambda^2/(Q^2(1-x))]^n$ 增长\cite[式~(5.37)]{Beneke1999}（原始分析见\cite{BenekeBraun1995b}；DIS 中泡链只产生 $1/Q^2$ 与 $1/Q^4$ 两种修正，见\cite{DasguptaWebber1996}）。需要强调的是，这一模型给出的是形状而非归一化，而且拟合数据所得的“扭度 4 修正”会随着扭度 2 项的微扰阶数提高而逐渐减小——二者在实践中不能被干净地分开\cite[第~5.2.4~节]{Beneke1999}，这正是\cref{conj:power-renormalon-ht}所描述的混合。
（以上 Mellin 关系都用符号计算核对过。）
\uses{prop:power-dis-twist4-x}
\end{example}
